\chapter{Derivational Morphology}\label{ch:derivational}

% ============================================================
% STATUS: SCAFFOLD — port and revise from
% archive/2020/chapters/06-derivational-morphology.tex.
% Key review points:
%   - Nominal derivation: update case terminology in examples
%   - Agent nominals: relation to the new antipassive stem (-kou)
%   - Abstract nominals: interaction with the new case system
%   - Verbal derivation: how does it interact with the antipassive /
%     applicative / causative slot? Are there derivational prefixes that
%     interact with valency?
%   - Borrowing: update to include recent borrowings and assimilation
%     patterns under the new stem-class system (the German -irná verbs)
% ============================================================

\section{Nominal derivation}

\subsection{Diminutives and augmentatives}

% [Port from 2020 archive]

\subsection{Agent nouns: \ird{-(h)ár} and \ird{-ist}}

Agent nouns denoting the habitual or professional performer of an action are formed with the suffixes \ird{-(h)ár} and \ird{-ist} (the latter a Slavic borrowing, now productive):

% [Examples: fisherman, teacher, activist, etc.]
% Note: the agent nominal -kou (from the antipassive stem) is a verbal
% derivation documented in §\,\ref{sec:nominalized}; the -ár/-ist forms
% are purely nominal derivation.

\subsection{Place nouns: \ird{-(i)mašt}}

% [Port from 2020 archive]

\subsection{Abstract nouns: \ird{-(i)žnám}}

% [Port from 2020 archive]

\section{Verbal derivation}

Unlike the valency operations described in Chapter~\ref{ch:verbs}, the
processes treated here do not occupy obligatory positions in the finite verbal
template. They form specialized lexical stems or stem variants, which then
take the ordinary inflectional morphology of the language. In the modern
language the most active stem-forming processes are denominal
verbalisation and, more marginally, reduplication.

\subsection{Denominal verbalisation}

Among overt noun-to-verb strategies, suffixation in \ird{-ěvat-} and
\ird{-it-} is the most productive. These are the predominant verbalising
formatives with learned, foreign, and recently borrowed nominal bases.
When a foreign stem is adapted directly as a verb, \ird{-ěvat-} is used
exclusively.

\pex
\a \ird{kontakt} $\rightarrow$ \ird{kontaktěvat-} \trsl{contact; get in touch}
\a \ird{telefon} $\rightarrow$ \ird{telefoněvat-} \trsl{telephone}
\xe

\subsubsection{Relational statives in \ird{-(e)zn-}}

A semantically specialized but productive denominal pattern adds
\ird{-(e)zn-} to a nominal base to form an intransitive relational stative
verb. The core meaning is \trsl{be X-like; belong to the domain of X; have
the character, style, material, or social colouring of X}. These verbs
conjugate exactly like the ordinary stative class described in
§\,\ref{sec:stative}, and they likewise admit the attributive derivative in
\ird{-á} and the manner adverb in \ird{-ě}
(§§\,\ref{sec:stative-attr}--\ref{sec:stative-adv}). Aspect follows the
usual stative pattern: the continuous gives plain state predication, the
progressive and perfective usually describe entry into the state, and the
retrospective describes a former state. Before this suffix, stem-final
\ird{k} surfaces as \ird{c}; thus \ird{barok} yields \ird{barocezna}.

Table~\ref{tab:relstatives} lists representative continuous forms together
with their derived attributives and manner adverbs.

\begin{table}[ht]
  \footnotesize\sffamily
  \caption{Relational stative verbs in \ird{-(e)zn-}.}\label{tab:relstatives}
  \medskip
  \begin{tabularx}{\textwidth}{L l l l L}
    \toprule
    {\sc base noun} & {\sc relational stative verb} & {\sc attributive} & {\sc manner adverb} & {\sc core sense} \\
    \midrule
    \ird{tvír} \trsl{court, courtly sphere} & \ird{tvírzna} & \ird{tvírzná} & \ird{tvírzně} & \trsl{be courtly; belong to the courtly register} \\
    \ird{ščech} \trsl{guild, craft corporation} & \ird{ščehezna} & \ird{ščehezná} & \ird{ščehezně} & \trsl{be guild-like; corporate, artisanal, guild-marked} \\
    \ird{mněra} \trsl{faith, creed} & \ird{mněrazna} & \ird{mněrazná} & \ird{mněrazně} & \trsl{be confessional; marked by a creed or religious orientation} \\
    \ird{kaměn} \trsl{stone} & \ird{kaměnzna} & \ird{kaměnzná} & \ird{kaměnzně} & \trsl{be stony; have a stone-like character; be stone-styled} \\
    \ird{sylo} \trsl{cloth, textile} & \ird{sylozna} & \ird{sylozná} & \ird{sylozně} & \trsl{be textile-like; be clothed, draped, or fabric-marked} \\
    \ird{barok} \trsl{baroque; baroque style} & \ird{barocezna} & \ird{barocezná} & \ird{barocezně} & \trsl{be baroque in style; belong to a baroque aesthetic} \\
    \ird{zmok} \trsl{temple, shrine} & \ird{zmocezna} & \ird{zmocezná} & \ird{zmocezně} & \trsl{be temple-like; be cultic, sacral, or shrine-associated} \\
    \ird{svíc} \trsl{peasant, smallholder, rustic estate} & \ird{svícezna} & \ird{svícezná} & \ird{svícezně} & \trsl{be peasantlike; be rustic or village-class marked} \\
    \ird{kmět} \trsl{law, order, code} & \ird{kmětezna} & \ird{kmětezná} & \ird{kmětezně} & \trsl{be juridical, code-bound, institutionalized} \\
    \bottomrule
  \end{tabularx}
\end{table}

\pex
\a \ird{tvírzná rěč} \trsl{courtly speech; speech of courtly register}
\a \ird{kaměnzná stěna} \trsl{a wall with a stony character} \lingcontext{a literal \trsl{stone wall} is expressed analytically, not by \ird{kaměnzná}}
\a \ird{barocezná dvorba} \trsl{a baroque-style building}
\xe

\subsection{Reduplication}\label{sec:reduplication}

Reduplication is a minor derivational strategy in modern Iridian. It is not an
exponent of grammatical aspect, does not mark alignment or valency, and is not
available as a fully general means of pluralizing events or participants.
Compared with languages in which reduplication is a central grammatical device,
its role in Iridian is sharply restricted. The construction belongs chiefly to
colloquial, expressive, and lightly lexicalized derivation.

The most useful synchronic description is that Iridian copies the initial
syllable of a base and prefixes that copy to the stem. In practice, however,
the pattern is only partly productive. Speakers accept it most readily with a
limited set of short verbal and adverbial bases, especially those denoting
repeated small actions, intermittent motion, tentative manipulation, or easily
segmentable manner. Many attested forms are therefore best treated as
conventionalized lexical items rather than as outputs of an unrestricted rule.

\pex
\a \ird{vad-} \trsl{think} $\rightarrow$ \ird{va-vad-} \trsl{ponder on and off; keep turning over in the mind}
\a \ird{ščen-} \trsl{arrive} $\rightarrow$ \ird{še-ščen-} \trsl{drop in from time to time; arrive intermittently}
\a \ird{skorně} \trsl{quickly} $\rightarrow$ \ird{sko-skorně} \trsl{quickly in short bursts; rather quickly}
\xe

Semantically, reduplication most often yields one of three closely related
readings. First, it may mark \term{attenuative iteration}: the event is
performed lightly, briefly, or in repeated small portions. Second, it may
contribute a \term{distributive} or \term{scattered} sense, especially with
motion and activity verbs, yielding meanings such as \trsl{here and there},
\trsl{from time to time}, or \trsl{in separate episodes}. Third, in a smaller
set of colloquial forms it serves an \term{expressive} or \term{tentative}
function, presenting the action as exploratory, hesitant, or half-committed.
Which of these readings is obtained depends chiefly on the lexical aspect of
the base and on the discourse context.

The contrast is lexical rather than inflectional. A reduplicated stem still
requires the ordinary aspect-alignment morphology of Chapter~\ref{ch:verbs},
and the resulting clause behaves morphosyntactically like any other finite
verb:

\pex
\a \ird{Janek list vadime.} \trsl{Janek is thinking about the letter.}
\a \ird{Janek list va-vadime.} \trsl{Janek is thinking about the letter on and off / mulling it over.}
\xe

For this reason reduplication should not be analyzed as part of the finite
template in §\,\ref{sec:morph-template}. It is better treated as a marginal
stem-forming operation whose output subsequently enters the ordinary verbal
complex. The reduplicated stem may in principle take the same voice, causative,
mood, and converb morphology as an unreduplicated stem, but in practice heavy
combinations are uncommon and often judged stylistically playful or forced.

Phonologically, reduplication introduces no special mutation class of its own.
Once the copied syllable has been prefixed, the derived form proceeds through
the regular morphophonemic system of Chapter~\ref{ch:phon}: weak-\ird{e}
deletion, grade formation where licensed, palatalization under front-vocalic
suffixes, and ordinary cluster repair apply exactly as they do elsewhere
(§\,\ref{sec:morphophon} and following). Apparent irregularities are usually
lexicalized remnants of older formations, not evidence for a separate
reduplicative phonology.

Outside the verbal domain, the language preserves a few related marginal
patterns of repetition, especially in estimative numeral expressions such as
\ird{měs a měs} \trsl{hundreds and hundreds / hundreds or so}. These are best
treated as peripheral analogues rather than as evidence that reduplication is a
major productive process in the grammar.

\section{Compounding}

% [Port from 2020 archive; update examples]

\section{Borrowing and assimilation patterns}

% The 2020 archive described:
%   - Phonological assimilation of loanwords (stress, consonant substitution)
%   - German loanwords in -irná (treated as a separate verbal class:
%     thematic consonant = -r, active root in -irž-)
%   - Recent borrowings (English, international) showing less assimilation
% These patterns remain valid in the 2026 framework; update terminology only.
